% Web専用原稿．
% 構成案：アルファベット・語・形式言語・集合演算と連接・Kleene閉包．
% 文献確認：所蔵文献検索で形式言語に特化した資料本文を確認できなかった．既出の有限列と自由モノイドから独自に展開する．
\section{形式言語}
% genron:begin K-000541
% genron:use K-000001
% genron:use K-000002
% genron:use K-000003
% genron:use K-000004
% genron:use K-000005
% genron:use K-000007
% genron:use K-000017
% genron:use K-000423
% genron:use K-000027
% genron:use K-000024
\begin{defn}[アルファベットと形式言語]
集合 $\Sigma$ をアルファベットといい，その自由モノイド $\Sigma^*$ の元を語という．形式言語とは $\Sigma^*$ の部分集合である．空語を $\varepsilon$ と書く．空言語 $\varnothing$ と空語のみを含む言語 $\{\varepsilon\}$ は相異なる．アルファベットは有限とは限らない．
\end{defn}
% genron:end K-000541
% genron:begin K-000542
% genron:use K-000001
% genron:use K-000002
% genron:use K-000003
% genron:use K-000004
% genron:use K-000005
% genron:use K-000007
% genron:use K-000017
% genron:use K-000541
% genron:use K-000423
% genron:use K-000118
% genron:use K-000108
\begin{prop}[語の長さ]
語 $u$ の有限列としての長さを $|u|$ とする．任意の語 $u,v$ に対し $|uv|=|u|+|v|$ である．また $|u|=0$ となる語は $\varepsilon$ に限る．
\end{prop}
\begin{proof}
連接は最初に $u$ の文字を，その後に $v$ の文字を置く有限列であるから長さは両者の長さの和である．長さ零の有限列は空列ただ一つである．
\end{proof}
% genron:end K-000542
% genron:begin K-000543
% genron:use K-000001
% genron:use K-000002
% genron:use K-000003
% genron:use K-000004
% genron:use K-000005
% genron:use K-000007
% genron:use K-000017
% genron:use K-000541
% genron:use K-000423
% genron:use K-000037
\begin{defn}[言語の連接]
形式言語 $L,M\subseteq\Sigma^*$ に対し $LM=\{uv\mid u\in L,\ v\in M\}$ をその連接と定める．和，共通部分，差及び補集合は母集合 $\Sigma^*$ における集合演算とする．
\end{defn}
% genron:end K-000543
% genron:begin K-000544
% genron:use K-000001
% genron:use K-000002
% genron:use K-000003
% genron:use K-000004
% genron:use K-000005
% genron:use K-000007
% genron:use K-000017
% genron:use K-000543
% genron:use K-000423
% genron:use K-000045
% genron:use K-000163
\begin{prop}[言語連接の演算法則]
言語の連接は結合律を満たし，$\{\varepsilon\}$ を単位元，$\varnothing$ を吸収元とする．任意の言語族 $(L_i)$ と言語 $M$ に対し
\[
 (\bigcup_i L_i)M=\bigcup_i(L_iM),\qquad
 M(\bigcup_iL_i)=\bigcup_i(ML_i)
\]
である．
\end{prop}
\begin{proof}
語の連接が結合的であるから $u(vw)=(uv)w$ であり，両集合の任意の元は相互に他方に属する．単位元と吸収元の性質は定義による．第一の分配律について，左辺の元はある $i$ の $u\in L_i$ と $v\in M$ の連接であり，右辺の元である．逆の包含も定義により成立する．第二も同様である．
\end{proof}
% genron:end K-000544
\section{反復と閉包}
% genron:begin K-000545
% genron:use K-000001
% genron:use K-000002
% genron:use K-000003
% genron:use K-000004
% genron:use K-000005
% genron:use K-000007
% genron:use K-000017
% genron:use K-000543
% genron:use K-000105
% genron:use K-000168
% genron:use K-000045
\begin{defn}[言語のKleene閉包]
言語 $L$ に対して $L^0=\{\varepsilon\}$，$L^{n+1}=L^nL$ と定め，
\[
 L^*=\bigcup_{n\in\mathbb N}L^n,\qquad
 L^+=\bigcup_{n\geq1}L^n
\]
と書く．$L^*$ を $L$ の Kleene 閉包という．
\end{defn}
% genron:end K-000545
% genron:begin K-000546
% genron:use K-000001
% genron:use K-000002
% genron:use K-000003
% genron:use K-000004
% genron:use K-000005
% genron:use K-000007
% genron:use K-000017
% genron:use K-000545
% genron:use K-000174
% genron:use K-000176
% genron:use K-000169
% genron:use K-000423
\begin{prop}[Kleene閉包の最小性]
$L^*$ は $L$ を含む自由モノイド $\Sigma^*$ の最小の部分モノイドである．従って $(L^*)^*=L^*$ である．
\end{prop}
\begin{proof}
空語は $L^0$ に属し，$L=L^1\subseteq L^*$ である．$u\in L^m$，$v\in L^n$ なら連接の結合律によって $uv\in L^{m+n}$ であるから積で閉じる．$L$ を含む部分モノイドは帰納法によって全ての $L^n$ を含むので $L^*$ を含む．最小性から冪等性も従う．
\end{proof}
% genron:end K-000546
% genron:begin K-000547
% genron:use K-000001
% genron:use K-000002
% genron:use K-000003
% genron:use K-000004
% genron:use K-000005
% genron:use K-000007
% genron:use K-000017
% genron:use K-000543
% genron:use K-000545
% genron:use K-000542
\begin{prop}[二文字の言語の二乗]
連接の二因子はともに $\varepsilon,a,b$ のいずれかである．空語を含む場合の結果は $\varepsilon,a,b$ であり，双方が一文字なら $aa,ab,ba,bb$ の四通りである．従って $L^2=\{\varepsilon,a,b,aa,ab,ba,bb\}$ である．
\end{prop}
% genron:end K-000547
